\documentclass[10pt,a4paper]{amsart}
\usepackage[margin=19mm]{geometry}
\usepackage{amsmath,amssymb,mathtools}
\usepackage[scr=rsfs,cal=euler]{mathalfa}
\usepackage{xfrac}
\usepackage[unicode,hidelinks,bookmarks=false]{hyperref}
\newcommand{\ie}{i.e.\ }
\newcommand{\cf}{cf.\ }
\newcommand{\resp}{respectively\ }
\newcommand{\Subsection}{\subsection}
\providecommand{\nobreakdash}{\nobreak\hbox{-}}
\setlength{\emergencystretch}{2em}
\multlinegap0pt
\usepackage{bm}
\usepackage{bbm}
\usepackage{dsfont}
\usepackage{xspace}
\usepackage{cancel}
\usepackage{mathtools}
\usepackage{amsthm}
\makeatletter
\@ifundefined{thm}{\newtheorem{thm}{Thm}}{}
\@ifundefined{lem}{\newtheorem{lem}[thm]{Lemma}}{}
\makeatother
\DeclareMathOperator{\im}{im}
\title[Formality of Sphere Bundles]{Formality of Sphere Bundles}
\author{Jiawei Zhou}
\date{Source: arXiv:2304.09594v3 (CC BY 4.0)}
\begin{document}
\maketitle

\noindent\emph{Source and licence.} Formality of Sphere Bundles. Version arXiv:2304.09594v3 (CC BY 4.0), distributed under the Creative Commons Attribution 4.0 International licence, as recorded in the arXiv licence metadata for this exact version and shown on the abstract page https://arxiv.org/abs/2304.09594v3. Full rights notice: licenses/arxiv-2304.09594-CC-BY-4.0.txt.\par\smallskip

\noindent\textit{Editorial scope.} Counted unit: the Lem whose source label is \texttt{vanishing uniform massey} in the article (source0.tex lines 244--250, source label \texttt{vanishing uniform massey}), delivered together with the article's own complete original proof of it (source0.tex lines 252--322, one \texttt{proof} environment of 71 lines). No mathematics is written, repaired or re-derived: statement and proof are the article's own text line for line. Required context retained uncounted: none. Every definition, display and label of those lines is retained unchanged. Omitted from this package and not redistributed: 1--243, 251--251, 323--608. Nothing inside a retained excerpt is omitted or altered: every retained body is delivered exactly as the article's lines are written, so no omission receipt of any kind accompanies this package. Citations: the retained excerpts carry no citation command at all, so no bibliography entry of the article is transcribed and no reference is redistributed. Reference adaptation: none was needed and none was made; every reference inside the delivered unit resolves inside it, so no unresolved reference or citation is produced. Printed slips kept literal: no printed word, symbol, hypothesis or number of the counted statement or of its proof was altered. Layout and class: the article's own class is \texttt{article} with packages including inputenc, amsmath, amsfonts, amssymb, amsthm, xy, enumerate, cite, authblk, hyperref; the wrapper is the lane's pinned \texttt{amsart} editorial wrapper at 10pt on A4, which restates the theorem-like environments the retained text needs and adds the article's own macro definitions that the retained text needs (\texttt{\textbackslash im}), with extra wrapper packages (bm, bbm, dsfont, xspace, cancel, mathtools). The delivered numbering is the unit's own. Assets: no retained span contains a figure environment, a graphics-inclusion command, a PDF-page inclusion, a table or a TikZ picture; no image file is copied into this package, the package declares no asset dependency and no image file is redistributed with it. Licence: version arXiv:2304.09594v3 of this article is distributed under the Creative Commons Attribution 4.0 International licence, as recorded in the arXiv licence metadata for this exact version and shown on the abstract page https://arxiv.org/abs/2304.09594v3. A full rights notice is delivered at licenses/arxiv-2304.09594-CC-BY-4.0.txt. Characters: every retained excerpt is pure ASCII, so the delivered engine, XeLaTeX with its default Latin Modern text font, reports no missing character.
\begin{lem}\label{vanishing uniform massey}
Let $\mathcal{A}$ be a CDGA with trivial differential. Suppose $\mathcal{A}_{\omega}$ is $n$-dimensional Poincar\'e and the Bianchi Massey tensor $\mathcal{F}:\mathcal{B}^{n+1}(\mathcal{A}_{\omega})\to H^n(\mathcal{A}_{\omega})$ vanishes. There are choices of $\alpha$ and $\gamma$ such that $\im\gamma\subset \mathbf{I}(\theta)$ and the morphism
$$
\gamma\alpha: E^*(\mathcal{A}_{\omega})\otimes H^*(\mathcal{A}_{\omega})\to \mathcal{A}_{\omega}^{*-1},\quad e\otimes a\mapsto \gamma(e)\alpha(a)
$$
is trivial on $K[E^*(\mathcal{A}_{\omega})\otimes H^*(\mathcal{A}_{\omega})]$.
\end{lem}

\begin{proof}
We will use $E^*$ to denote $E^*(\mathcal{A}_{\omega})$ in this proof. It is sufficient to construct $\alpha$ and $\gamma$ such that $\im\gamma\subset \mathbf{I}(\theta)$ and the uniform Massey product $\mathcal{T}=0$ on $K[E^*\otimes H^*(\mathcal{A}_{\omega})]$. Then the image of $\gamma\alpha$ are exact elements in $\mathbf{I}(\theta)$. But as discussed earlier in this section, the exact elements in $\mathcal{A}_{\omega}$ are in $\im\omega$. So $\gamma\alpha$ has to be 0 on $K[E^*\otimes H^*(\mathcal{A}_{\omega})]$.

Pick a right inverse of the left multiplication map $\omega:\mathcal{A}\to\im\omega$ and extend it to an endomorphism on $\mathcal{A}$. We use $\omega^{-1}$ to denote this map.

Choose an $\alpha:H^*(\mathcal{A}_{\omega})\to \mathcal{A}_{\omega}^*$ such that $\alpha(\theta\ker\omega)\subset\mathbf{I}(\theta)$. As $\alpha^2(E^*)$ consists of exact forms, it is in $\im\omega\subset\mathcal{A}$. So we can set
$$
\gamma=\theta\circ\omega^{-1}\circ\alpha^2:E^*\to \mathcal{A}_{\omega}^{*-1}.
$$
Then $\im\gamma\subset \mathbf{I}(\theta)$. It follows that $\im\mathcal{T}\subset\theta\ker\omega$. For an extreme case that $\ker\omega=0$, we already have $\mathcal{T}=0$. Note that in this case $\mathcal{A}$ is infinite dimensional, unless $\mathcal{A}=0$.

From now on we consider the general case that $\theta\ker\omega$ is non-trivial. Then $\theta\ker\omega$ must contain $H^n(\mathcal{A}_{\omega})$, because for non-trivial $\theta x\in\theta\ker\omega$ there exists $y\in H^*(\mathcal{A}_{\omega})$ such that
$$
\alpha_H(\theta xy)=(\alpha_H\frown \theta x)y\neq 0.
$$

To make $\mathcal{T}=0$, we will change $\gamma$ to $\gamma'=\gamma+\eta$ for some $\eta:E^*\to \mathcal{Z}^{*-1}(\mathcal{A}_{\omega})$ satisfying $\im\eta\subset \mathbf{I}(\theta)$.

Consider the map
$$
\gamma\alpha^2:E^*\otimes H^*(\mathcal{A}_{\omega}) \otimes H^*(\mathcal{A}_{\omega})\to \mathcal{A}_{\omega}^{*-1},\quad e\otimes x\otimes y\mapsto \gamma(e)\alpha(x)\alpha(y).
$$
It takes closed values on $K[E^*\otimes H^*(\mathcal{A}_{\omega}) \otimes H^*(\mathcal{A}_{\omega})]$, and factors through the projection $E^*\otimes H^*(\mathcal{A}_{\omega}) \otimes H^*(\mathcal{A}_{\omega})\to E^*\otimes \mathcal{G}^2H^*(\mathcal{A}_{\omega})$. Hence, $\gamma\alpha^2$ induces another map
$$
\mu:K[E^*\otimes \mathcal{G}^2H^*(\mathcal{A}_{\omega})]\to H^{*-1}(\mathcal{A}_{\omega}).
$$
$\mu$ depends on the choice of $\gamma$. Our goal is to find some $\gamma'$ such that the corresponding $\mu'=0$ when acting on the degree $n+1$ part of $K[E^*\otimes \mathcal{G}^2H^*(\mathcal{A}_{\omega})]$.

First consider the restriction of $\mu$ to $K[E^*\otimes E^*]$. Note that here $K[E^*\otimes E^*]$ means the kernel of full graded symmetrisation $E^*\otimes E^*\to \mathcal{G}^4H^*(\mathcal{A}_{\omega})$. For arbitrary $e,e'\in E^*$,
\begin{align*}
\gamma(e)\alpha^2(e')-(-1)^{|e||e'|}\gamma(e')\alpha^2(e) &= \gamma(e)d\gamma(e')-(-1)^{|e||e'|+(|e'|-1)|e|}d\gamma(e)\gamma(e') \\
&= (-1)^{|e|-1}d\big(\gamma(e)\gamma(e')\big).
\end{align*}
So $\mu$ vanishes on graded anti-symmetric tensors, and it factors through the projection $K[E^*\otimes E^*]\to\mathcal{B}^*(\mathcal{A_{\omega}})$. Moreover, the induced map $\mathcal{B}^*(\mathcal{A_{\omega}})\to H^{*-1}(\mathcal{A}_{\omega})$ is exactly the Bianchi-Massey tensor, which is 0 when acting on $\mathcal{B}^{n+1}(\mathcal{A_{\omega}})$ by assumption. Therefore, $\mu\equiv 0$ on $K[E^*\otimes E^*]$, in particular on its degree $n+1$ part. As $\eta(e)\alpha^2(e')$ is exact for any $e,e'\in E^*$, $\mu'=0$ on the degree $n+1$ part of $K[E^*\otimes E^*]$ for any choice of $\gamma'$.

Let $q:\mathcal{G}^2H^*(\mathcal{A}_{\omega})\to \mathcal{G}^2(\text{coker}\,\omega)$ be the morphism induced by the projection $H^*(\mathcal{A}_{\omega})\to \text{coker}\,\omega$ with kernel $\theta\ker\omega$, and $L^*$ be the preimage of $\theta\ker\omega$ under the multiplication map $c:\mathcal{G}^2H^*(\mathcal{A}_{\omega})\to H^*(\mathcal{A}_{\omega})$. Then $c\circ q(L^*)=0$ because its image is the projection of $c(L^*)$. Hence, $q(L^*)\subset E^*$. So there exists an induced map
$$
Q: E^*\otimes L^*\to E^*\otimes E^*,\quad e\otimes a\mapsto qe\otimes qa.
$$
Moreover, $Q$ sends $K[E^*\otimes L^*]$ into $K[E^*\otimes E^*]$. Here $K[E^*\otimes L^*]$ also means the kernel of full graded symmetrisation $E^*\otimes L^*\to \mathcal{G}^4H^*(\mathcal{A}_{\omega})$.

Let $e\in E^*$ and $a\in L^*$ with $|e|+|a|=n+1$. Then every term of $a-qa$ has a factor in $\theta\ker\omega$. Thus, $\alpha^2(a-qa)\in\mathbf{I}(\theta)$. As $\gamma(e)$ is also in $\mathbf{I}(\theta)$, we have $\gamma(e)\alpha^2(a)=\gamma(e)\alpha^2(qa)$. For the same reason,
\begin{align*}
\gamma(e)\alpha^2(qa) &= (-1)^{|e|}d\big(\gamma(e)\gamma(qa)\big)-(-1)^{|e|}\alpha^2(e)\gamma(qa) \\
&= (-1)^{|e|}d\big(\gamma(e)\gamma(qa)\big)-(-1)^{|e|}\alpha^2(qe)\gamma(qa) \\
&= (-1)^{|e|}d\big(\gamma(e)\gamma(qa)\big)-(-1)^{|e|}d\big(\gamma(qe)\gamma(qa)\big)+\gamma(qe)\alpha^2(qa).
\end{align*}
Therefore, $\mu(e\otimes a)=[\gamma(e)\alpha^2(a)]=[\gamma(qe)\alpha^2(qa)]=\mu\circ Q(e\otimes a)$. As discussed above, $\mu$ vanishes on $K[E^*\otimes E^*]$. So $\mu(K[E^*\otimes L^*])=\mu\circ Q(K[E^*\otimes L^*])\subset\mu(K[E^*\otimes E^*])=0$.

Now let $p=id\otimes c: K[E^*\otimes\mathcal{G}^2H^*(\mathcal{A}_{\omega})]\to E^*\otimes H^*(\mathcal{A}_{\omega})$ sending $e\otimes(x\cdot y)$ to $e\otimes xy$. Then $\ker p$ is exactly $K[E^*\otimes E^*]$. Hence, $p$ induces a morphism $\bar{\mu}:\im p\to H^*(\mathcal{A}_{\omega})$ of degree $-1$. As $\mu$ vanishes on $K[E^*\otimes L^*]$, $\bar{\mu}=0$ on $\im p\cap (E\otimes\theta\ker\omega)$. So we can extend $\bar{\mu}$ to all $E^*\otimes H^*(\mathcal{A}_{\omega})$ such that it vanishes on $E^*\otimes\theta\ker\omega$.

For each $e\in E^i$, $\bar{\mu}$ induces a morphism
$$
H^{n+1-i}(\mathcal{A_{\omega}})\to \mathbb{K}, \quad x\mapsto -\alpha_H\circ\bar{\mu}(e\otimes x)
$$
where $\alpha_H$ is the Poincar\'e class of $H^*(\mathcal{A_{\omega}})$. Let $\delta(e)$ be the unique element in $H^{i-1}(\mathcal{A_{\omega}})$ such that this morphism equals to $\alpha_H\frown\delta(e)$. We claim that $\delta(e)\in\theta\ker\omega$.

Write $\delta(e)=x_0+\theta y_0$ for some $x_0\in\text{coker}\,\omega$ and $y_0\in\ker\omega$. If $x_0\neq 0$, there exist some $x_1\in\text{coker}\,\omega$ and $y_1\in\ker\omega$ such that $(\alpha_H\frown x_0)(x_1+\theta y_1)\neq 0$, i.e. $x_0(x_1+\theta y_1)$ is non-trivial in $H^n(\mathcal{A_{\omega}})$. By the discussion at the beginning of the proof, $H^n(\mathcal{A_{\omega}})\subset \theta\ker\omega$. So $x_0x_1$ has to be 0. It follows that $\delta(e)\theta y_1=(x_0+\theta y_0)\theta y_1=x_0\theta y_1\neq 0$. On the other hand, $\delta(e)\theta y_1=-\alpha_H\circ\bar{\mu}(e\otimes\theta y_1)=0$ as $\bar{\mu}=0$ on $E^*\otimes\theta\ker\omega$, which is a contradiction. Hence, $x_0$ must be 0 and $\delta(e)\in\theta\ker\omega$.

Let $\eta=\alpha\circ\delta:E^*\to \mathcal{Z}^{*-1}(\mathcal{A_{\omega}})$. It represents $\delta$ and its image is in $\mathbf{I}(\theta)$. Then $\gamma'\alpha^2=(\gamma+\eta)\alpha^2$ induces a map $\mu':K[E^*\otimes \mathcal{G}^2H^*(\mathcal{A}_{\omega})]\to H^{*-1}(\mathcal{A}_{\omega})$. For arbitrary $\sum e_j\otimes a_j\in K[E^*\otimes \mathcal{G}^2H^*(\mathcal{A}_{\omega})]$,
\begin{align*}
\mu'\left(\sum e_j\otimes a_j\right) &= \mu\left(\sum e_j\otimes a_j\right)+\sum [\eta(e_j)\alpha^2(a_j)] \\
&= \bar{\mu}\left(\sum e_j\otimes c(a_j)\right)+\sum\delta(e_j)c(a_j) = 0.
\end{align*}

Finally, we prove that the uniform Massey product $\mathcal{T}'$ defined by $\gamma'$ is identically 0. This follows from $K[E^*\otimes H^*(\mathcal{A}_{\omega})]\otimes H^*(\mathcal{A}_{\omega})\subset K[E^*\otimes H^*(\mathcal{A}_{\omega})\otimes H^*(\mathcal{A}_{\omega})]$ and
$$
\mathcal{T'}\left(\sum e_j\otimes x_j\right) y=\left[\sum\gamma(e_j)\alpha(x_j)\alpha(y)\right]=\mu'\left(\sum e_j\otimes (x_j\cdot y)\right)=0
$$
for $e_j\in E^*$, $x_j,y\in H^*(\mathcal{A}_{\omega})$ and $|e_j|+|x_j|+|y|=n+1$. Thus, $\alpha_H\frown \mathcal{T'}\left(\sum e_j\otimes x_j\right)=0$ and then $\mathcal{T'}\left(\sum e_j\otimes x_j\right)=0$. As $\im\gamma'\subset\mathbf{I}(\theta)$ by construction, the discussion in the first paragraph of the proof shows that $\gamma'\alpha$ has to be 0 on $K[E^*\otimes H^*(\mathcal{A}_{\omega})]$. 
\end{proof}

\end{document}
